\documentclass[a4paper, 12pt]{book}
\usepackage{amsmath, amssymb, amsfonts, amsthm}
\usepackage{physics}
\usepackage{geometry}
\usepackage{hyperref}
\newtheorem{definition}{Definition}[chapter]
\newtheorem{theorem}{Theorem}[chapter]
\newtheorem{lemma}{Lemma}[chapter]
\newtheorem{example}{Example}[chapter]
\newtheorem{proposition}{Proposition}[chapter]
\setcounter{secnumdepth}{3}
\title{Notes for MAT267}
\author{Eric Chang}
\date{Spring 2026}
\begin{document}
\maketitle
\tableofcontents
\chapter{Solving Simple ODEs}
\chapter{Existence and Uniqueness of Solutions}
\section{Introduction}
In this chapter, our biggest goal is to investigate the existence and uniqueness of solutions to an ODE.
\section{Basic Notions}
\subsection{Continuity}
Consider the setting \( E \subset \mathbb{R}^n \) and \( f: E \to \mathbb{R}^m \)
\begin{definition}
	We say that \( f \) is continuous at a point \( a \in E \) if it satisfies: \[
		\forall \varepsilon>0\quad \exists \delta>0 \quad \forall x \in E \quad (\abs{x-a} < \delta \implies \abs{f(x)-f(a)}< \epsilon)	    \]
\end{definition}
\begin{definition}
	Alternatively, we say a function \( f:E \to \mathbb{R}^m \) is uniformly continuous on a set \( D \subset E \) \[
		\forall \varepsilon>0 \quad \exists \delta > 0 \quad \forall x,y \in D\quad(\abs{x-y}<\delta\implies \abs{f(x)-f(y)}<\varepsilon)
	\]
\end{definition}
See that we can say a function is continuous on \( D \subset E \) if it is continuous in each point on \( D \), but that doesn't necessarity satisfy uniform continuity on \( E \) as they might not be a "global" bound on the domain such that the \( \varepsilon \) bounds the range.\\

We can take this "global" bound concept further to a family of functions:
\begin{definition}
	We call a possibly infinite family of function \( \mathcal{F}\subset \{f:E\to \mathbb{R}^m \}\) \textbf{equicontinuous} on \( D \subset \mathbb{R}^m \)if \[
		\forall \varepsilon>0 \quad \exists \delta>0\quad \forall f \in \mathcal{F}\quad \forall x,y\in E \quad (\abs{x-y} <\delta \implies \abs{f(x)-f(y)} < \varepsilon)
	\] for any \( x,y \in D \).
\end{definition}
In other words, a set of functions that is equicontinous on \( D \) must have functions that are individually uniformly continuous, and also for each \( \varepsilon \), the set of \( \delta \) satisfing \( \varepsilon \) for \( f\in \mathcal{F} \) must have a non-zero infimum.
\subsubsection{Examples}
a
\subsection{Uniform Convergence}
We define a norm on functions:
\begin{definition}
	The Uniform Norm is defined for a function \( f:E\to \mathbb{R}^m \):
	\[
		\norm{f} _\infty =\operatorname{sup}_{x\in E} f(x)
	\]
\end{definition}
and a type of convergence of a sequence of functions
\begin{definition}
	For a sequence of functions \( f_n: E\to \mathbb{R}^m \), we say that the sequence of functions converges uniformly to \( f:E\to \mathbb{R}^m \) if
	\[
		\forall \varepsilon>0 \quad \exists N >0 \quad \forall n \in\mathbb{N} \quad \forall x \in E \quad (n>N \implies \abs{f_n(x)-f(x)}<\varepsilon)
	\]
	We denote this as \[
		f_n \rightrightarrows f
	\]
\end{definition}
This can be easily interpreted as a improvement of the natural form of convergence, which provides possibly varying \( \delta_x \) to satisfy \( \varepsilon \) for each \( x\in E \). There are many implications of Uniform convergence of a sequence of function.
\begin{theorem}
	If \( f_n, f:E\to \mathbb{R}^m \) and each \( f_n \) is continuous on \( E \), then \( f \) is continuous.
\end{theorem}
\begin{proof}
	\[
		\norm{f(x)-f(a)}\leq \norm{f_n(x)-f(x)} -\norm{f_n(x)-f_n(a)} -\norm{f_n(a)-f(a)}< 3 \varepsilon
	\]
	by choosing \( n \) and \( \delta \).
\end{proof}
\begin{theorem}
	Let \( E \) be a compact subset of \( \mathbb{R}^n \), \( f_n\in C(E,\mathbb{R}^m) \). If \( f_n \rightrightarrows f \) for some \( f \), then
	\begin{enumerate}
		\item \( \{f_n\} \) is equicontinuous
		\item \( \{f_n\} \) is uniformly bounded in the uniform norm, that is \[
			      \exists M\in \mathbb{R}^+ \quad \forall n \in \mathbb{N} \quad (\norm{f_n}_\infty \leq M)
		      \]
	\end{enumerate}
\end{theorem}
\subsection{Ascoli-Arzela}
\begin{theorem}
	If \( E \) is a compact subset of \( \mathbb{R}^n \) and some infinite family \( \mathcal{F} \subset C(E,\mathbb{R}^m) \). If \( \mathcal{F} \) is equicontinous and unformly bounded, then for every sequence in \( \mathcal{F} \) there is a convergent subsequence.
\end{theorem}
\begin{proof}
	Not required.
\end{proof}
\subsection{Relative Compactness}
\begin{definition}
	Let \( K \) be a compact subset of \( \mathbb{R}^n \). A set \( A \subset C(K, \mathbb{R}^m) \) is called \textbf{relatively compact} if \( \overline{A} \) is compact. That is to be said, if \( A \) is a metric space, every sequence in \( A \) has a convergent subsequence.
\end{definition}
\begin{theorem}
	If \( K \) is a compact subset of \( \mathbb{R}^n \) and \( A\subseteq C(K, \mathbb{R}^m) \), then \( A \) is relatively compact if and only if
	\begin{enumerate}
		\item \( A \) is equicontinuous
		\item \( A \) is uniformly bounded
	\end{enumerate}
\end{theorem}
\begin{proof}
	\( \implies \)\\
	If \( A \) is relatively compact, then we know that every sequence has a convergent subsequence under the unform norm. In addition, \( \overline{A} \) is compact, thus \( A \) is bounded by the uniform norm. If \( \mathcal{F} \) is not equicontinuous, then for some \( \varepsilon>0 \) and any selection of \( \delta>0 \) there is a \( f \in \mathcal{F}\) such that \( \norm{f(x)-f(y)} <\varepsilon \) is not satisfied by \( \norm{x-y} <\delta \). \[
		\norm{f(x)-f(y)} \geq \varepsilon
	\]\\
	But recall that every sequence has a converging subsequence. Consider an converging sequence of such \( f \) \[
		f_k \rightrightarrows g
	\]
	We know that \( g \) must be uniformly continuous as \( g \) is continuous on a compact set. Thus, for all \( \frac{ 1 }{ 3 } \varepsilon>0 \) there must exist a \( \delta>0 \) such that \( \norm{x-y} <\delta \)\[
		\norm{g(x)-g(y)} <\frac{ \varepsilon }{ 3 }
	\]
	Since \( f_k \) converges uniformly to \( g \), there must be a \( N\in \mathbb{N}\) such that \( k>m \) means that \[
		\norm{f_k(x)-g(x)} <\frac{ \varepsilon }{ 3 }
	\]
	Then, \[
		\norm{f_k(x)-f_k(y)}<\varepsilon
	\]
	for \( k>N \) and \( \norm{x-y} <\delta \) which contraducts the assumption that \( \norm{f(x)-f(x)}\geq \varepsilon  \) for some \( \varepsilon, x,y \) where \( \norm{x-y} <\delta \) for any \( \delta \).
	\\
	\(\Leftarrow\)
	\\
	If equicontinuous and uniformly bounded, then by Ascoli-Arzela every sequence has a convergent subsequence, thus \( A \) is relatively compact by definition.
\end{proof}
\subsubsection{Radial Extention}
\begin{lemma}
	Let \( f: \overline{\mathbb{B}}(0,R)\subset \mathbb{R}^n \to \mathbb{R} \) be continuous, then we can extend the ball to the entire domain \( \mathbb{R}^n \) by defining \( F:\mathbb{R}^n\to \mathbb{R}^m \) by \[
		F(x) =
		\begin{cases}
			f(x),                       & \norm{x} \leq R \\
			f(\frac{ x }{ \norm{x}  }), & \norm{x} >R
		\end{cases}
	\]
	Then \( F \) is continuous on \( \mathbb{R}^n \) and extends \( f \).
\end{lemma}
\section{Fixed Point Theorems}
\subsection{Banach's Contraction Mapping Theorem}
\begin{theorem}[Banach's Contraction Mapping theorem]
	Let \( (D,d) \) be a complete metric space and \( T:D\to D \) be a string contraction, that is, \[
		\exists q\in (0,1) \quad \forall x,y\in D \quad (d(T(x),T(y))\leq d(x,y)
	\]
	Then \( T \) has a \textbf{unique} fixed point in \( D \).
\end{theorem}
\begin{proof}
	We use Picard's Iteration to solve the problem. Define the iteration
	\[
		x(n)=T(x_{n-1} )
	\]
	and let \( x_0\in D \) be arbitary. We first want to show that this sequence converges to some \( x^* \in D \).
	See that \[
		d(x_{n+1} , x_n)\leq d(T(x_n), T(x_{n-1} ))\leq q d(x_n,x_{n-1})\leq q^nd(x_1,x_0)
	\]
	thus we have that \[
		d(x_n,x_m)\leq \sum_{i=n}^{m-1} d(x_{i+1}, x_i)\leq d(x_1,x_0)\sum_{i=n}^{m-1} q^i
	\]

	Since \[
		\sum_{i=n}^{m} q^i\leq \sum_{i=1}^{\infty} q^i-\sum_{i=1}^{n} q^i= \frac{ 1 }{ 1-q } - S_n
	\]
	Since \( \lim_{n \to \infty} S_n=\frac{ 1 }{ 1-q }  \), we have that \[
		\forall \varepsilon>0 \quad \exists N>0 \quad \forall n \in \mathbb{N} \quad (n>N \implies \sum_{i=n}^{m} q^i \leq \frac{ 1 }{ 1-q } -S_n <\varepsilon)
	\]
	Thus \( x_n \) is cauchy, and as \( D \) is complete, it must converge to some \( x^* \) in \( D \). We verify that \( x^* \) is a fixed point: \[
		d(T(x^*), x^*)\leq d(T(x^*),T(x_k))+d(x_k, x^*)\leq (q+1)d(x^*,x_k)<\varepsilon
	\]
	Where \( \varepsilon \) can be made arbitarily small by selection of \( k \) by convergence. Thus, \[
		T(x^*)=x^*
	\]
	We now want to show uniqueness. This can be done the same way:
	\[
		d(x^*,y^*)=d(T(x^*),T(y^*))\leq qd(x^*,y^*)
	\]
	Thus \[
		d(x^*,y^*)=0
	\]
	Then \( x^*=y^* \)
\end{proof}
\subsection{Schauder-Tychonoff Theorem}
We begin by constructing a necessary tool:
\begin{theorem}
	[Brouwer's Fixed Point Theorem]
	Let \( B=\overline{B}(0,1) \) be a closed unit ball of \( \mathbb{R}^n \) and \( T:B\to B \) be continuous. Then \( T \) has atleast one fixed point in \( B \).
\end{theorem}
\begin{proof}
	The proof is not required.
\end{proof}
Now we proceed to the actuall theorem
\begin{theorem}
	[Schauder-Tychonoff Fixed Point Theorem]
	Let \( E=C([0,1], \mathbb{R}^n) \) and \( (E,\norm{\cdot})\) a Banach Space.  Let \( B \) be the unit ball of \( E \): \[
		B=\{ f\in E:  \norm{f} \leq 1\}
	\]
	Assume that \( T:B\to B \) is continuous with respect to the uniform norm and \( T(B) \) relatively compact in \( E \), then \( T \) has a fixed point in \( B \).
\end{theorem}
\begin{proof}
	We wish to use picard's iteration again: construct a suitable sequence that converges to the fixed point. We construct this sequence by linear interpolation:\\

	For each \( k\in \mathbb{N} \), pick \( a_0, a_1, \cdots , a_k \in \overline{B}(0,1)\subset \mathbb{R}^n\). Define a function \[
		x_k:[0,1]\to \mathbb{R}^n
	\]
	such that for \( j\in \mathbb{N} \) and \( j\leq k \), \( t\in[\frac{ j-1 }{ k } , \frac{ j }{ k } ] \):
	\[
		x_k(t)=a_{j-1} +(a_j-a_{j-1})(t-\frac{ j-1 }{ k } )
	\]
	then \( x_k(\frac{ i }{ k }) = a_i \) and \( x_k \) is continuous. And since \( ||a_k||\leq 1 \) inside the closed ball of radius \( 1 \), and \( x_k(t) \) consists of points defined on lines between \( a_k \)'s, \( x_k([0,1])\) is bounded by \( \overline{B}(0,1)\subset \mathbb{R}^n \), thus \( x_k\in B \).
	\\
	Let \( B_k \) denote the set of all \( x_k \) (with arbitary selection of \( a_j \)), with the map:
	\[
		P_k(a_0, \cdots, a_k)= x_k \in B_k, \quad x_k(\frac{ i }{ k }) =a_i
	\]
	\( P_k \) is continuous as a small change in \( a_i \) corresponds to a small change in \( x_k \) due to stability of linear interpolation. Now let
	\[
		T_k(a_0,\cdots, a_k)=(T(P_k (a_0,\cdots, a_k)(\frac{ 0 }{ k } ), \cdots , T(P_k(a_0, \cdots, a_k))(\frac{ k }{ k } ))
	\]
	Then \( T_k \) is continuous by investigating each component, with \[
		||T\circ P_k(a)(\frac{ i }{ k }) - T\circ P_k(b)(\frac{ i }{ k } )||\leq \norm{T\circ P_k(a)-T\circ P_k (b)} _\infty
	\]
	where continuity follows from continuity under uniform norm.
	Thus, \( T_k: \overline{B}(0,1) ^{k+1} \to \overline{B}(0,1)^{k+1}  \), and is continuious, where \( \overline{B}(0,1)^k \) is homoemorphic to \( \overline{B}(0,1) \subset \mathbb{R}^{n(k+1)}  \). Thus, apply Brouwer's Fixed Point Theorem. Then, there is a fixed point \( a^*=(a^*_0, \cdots, a^*_k) \) of \( T_k \), where \[
		T_k(a^*)=(a^*_0,\cdots, a^*_k), \quad T(P_k(a^*))(\frac{ i }{ k } )=P_k(a^*)(\frac{ i }{ k } )=a^*_i
	\]
	Then we have a partially fixed point \( x^*_k=P^k(a^*) \), for which \( Tx^*_k=x^*_k \) for the interpolated points $a^*$. Construct a sequence of such \( (T(x^*_k))_{k\in \mathbb{N}} \). See that since \( T(x_k^*)\in T(B)\) is relatively compact, there is a convergent subsequence converging by the uniform norm to some \( x^* \in B\). Index this sequence by \( x_k \). \\

	We want to show that \( \norm{T(x_k)-x_k}_\infty \rightarrow 0\). That is, we need to show that \[
		\forall \varepsilon>0 \quad \exists N>0 \quad \forall x \in [0,1] \quad (k>N \implies \norm{Tx_k-x_k}_\infty<\varepsilon )
	\]
	See that \[
		\norm{Tx_k-x_k} =\sup_{t\in[0,1]}\norm{T(x_k)(t)-x_k(t)}
	\]
	Then by triangle inequality for \( t\in[\frac{ i-1 }{ k } , \frac{ i }{ k } ] \):
	\[
		\norm{Tx_k(t)-x_k(t)} \leq \norm{Tx_k(t)-Tx_k(\frac{ i }{ k })} + \norm{Tx_k(\frac{ i }{ k })-x_k(\frac{ i }{ k } ) } +\norm{x_k(\frac{ i }{ k } )-x_k(t)}
	\]
	Let \( \varepsilon>0 \). Since \( T(B) \) is relatively compact, it must be equicontinuous. Thus, there is a \( delta>0 \) such that \( \norm{x-y} <\delta \implies \norm{Tx_k(x)-Tx_k(y)}<\frac{ \varepsilon }{ 2 }\). Then choose \( N\) large enough such that \( \frac{ 1 }{ k } <\delta \). \\
	Also see that \( ||x_k(\frac{ i }{ k }-x_k(t)||\leq \norm{x_k(\frac{ i }{ k } )-x_k(\frac{ i-1 }{ k } )}=\norm{Tx_k(\frac{ i }{ k } )- Tx_k(\frac{ i-1 }{ k } )} \).
	\\

	Then, if \( k>N \), \[
		\norm{Tx_k(t)-x_k(t)}<\varepsilon
	\]
	for all \( t\in [\frac{ i-1}{k } , \frac{ i }{ k } ] \), and \( N \) independent of \( i \), this is true for all \( t\in[0,1] \). Then, \[
		x_k \rightrightarrows x^*
	\]
	To verify that this is indeed a fixed point, \[
		T(x^*)=T(\lim_{k \to \infty } x_k)=\lim_{k \to \infty}T(x_k)= x^*
	\]
	And thus \( x^*\) is a fixed point in \( B \)

\end{proof}
\section{Existence of Solutions}
\begin{definition}
	[The Global Cauchy Problem]
	Let \( t_0<t_1 \) in \( \mathbb{R} \) and \[
		f:[t_0,t_1]\times \mathbb{R}^n \to \mathbb{R}^n
	\]
	together with an inital value \( x_0\in \mathbb{R}^n \). The initial value problem is :
	\[
		\begin{cases}
			x^\prime(t)=f(t,x(t)), & t\in[t_0,t_1] \\
			x(t_0)=x_)
		\end{cases}
	\]
	where we seek a function \( x\in C^1([t_0,t_1], \mathbb{R}^n \) if possible.
\end{definition}
\begin{definition}
	[Intergral Equation for an IVP]
	Given \( f \) and \( x_0 \) as above, the intergral equation is given by \[
		x(t)=x_0+\int_{t_0}^{t} f(s,x(s)) \dd s
	\]
\end{definition}
A solution solves the IVP if and only if it satisfies the Integral Equation.
\\

With this as motivation, we define the picard map :
\begin{definition}
	For \( (f,x_0,t_0,t_1) \) as previously defined, the picard map for the IVP is given by:
	\[
		T(x)=x_0+\int_{t_0}^{t} f(s,x(s)) \dd s
	\]
	Then the solution to the IVP is exactly the fixed points of \( T \).
\end{definition}
and with these definitions, we have our first theorem for existence of solution
\subsection{Cauchy-Peano existence Theorem}
\begin{theorem}
	[Cauchy-Peano Existence]
	Consider a global cauchy problem as defined with \( (f,x_0,t_0,t_1) \) . If \( f \) is continuous and  bounded, then the IVP has at least one solution on \( [t_0, t_1]) \)\end{theorem}
There are two main ways to prove this:
\subsubsection{Proof 1}
\begin{proof}
	Consider Banach Space \( C([t_0,t_1]) \) with the supremum norm. Let \[
		B=\{x\in C([t_0,t_1], \mathbb{R}^n): \norm{x} _\infty \leq 1\}
	\]
	Define the Picard Map:
	\[
		Tx=x_0+\int_{t_0}^{t_1} f(s,x(s)) \dd s
	\]
	We wish to find a fixed point of \( T \) on \( B \). The idea is to use the Schauder-Tychonoff theorem. Thus we need to show that \( T \) is continuous and \( T(B) \) is relatively compact. \\

	To show that \( T \) is continuous. It is sufficient to show that any uniformly convergent sequence \( (x_k)\rightrightarrows x \) means that \( (T(x_k))
	\rightrightarrows T(x)\), that is if \( x_k \to x \), then
	\[
		||T(x_k)-T(x)||_\infty=\sup_{t\in[t_0, t_1]} \norm{T(x_k)(t)-T(x)(t)}
	\]
	Thus,
	\begin{align*}
		\norm{T(x_k)(t)-T(x)(t)} & \leq \norm{\int_{t_0}^{t} f(s,x_k(s))-f(s,x(s))\dd s}         \\
		                         & \leq (t_1-t_0) \sup_{s\in[t_0,t_1]} ||f(s,x_k(s))-f(s,x(s))||
	\end{align*}
	but since \( f \) is continuous at \( x \), there exists \( N>0 \) such that \( k>N \) means that \( \norm{f(x_k)-f(x)} <\epsilon \). Then \( \norm{T(x_k)-T(x)} _\infty < \varepsilon \) up to choice of \( N \), thus it is continuous. \\
	Now we want to show that it \( T(B) \) is relatively compact, that is , \( T(B) \) is uniformly bounded and equicontinuous. See that
	\begin{align*}
		||T(x)(t)|| & = ||x_0|| + \int_{t_0}^{t} ||f(s, x(s))||\dd s \\
		            & \leq ||x_0|| +M(t_1-t_0)
	\end{align*}
	Thus is uniformly bounded. With \[
		H=\norm{x_0} +M(t_1-t_0)
	\] To show equicontinuity, consider
	\begin{align*}
		\norm{T(x)(t)-T(x)(t^\prime)} & =\int_{t^\prime}^{t} \norm{f(s,x(s))} \dd s \\
		                              & \leq M|t-t^\prime|
	\end{align*}
	Then \( T \) is equicontinuous since the bound is independent of \( x \).
	Now consider two cases:
	\\
	\textbf{\( H\leq 1 \)}\\
	Then we may directly apply Schauder-Tychonoff as \( T:B \to B \). Since \( T \) is equicontinuous and uniformly bounded, it is relativaly compact. Thus, by Schauder -Tychonoff, there is a fixed point \( x^*:[t_0,t_1]\to \mathbb{R}^n \) a solution to the IVP. \\
	\textbf{\( H>1 \)}\\
	Then we scale :
	\[
		y_0=x_0 \quad g(t,y(t))=\frac{1}{H} f(t, H x(t)), \quad \text{set} x(t)=Hy(t)
	\]
	Then Schauder-Tychonoff follows similarly.
\end{proof}
\subsubsection{Proof 2}
\subsection{Picard-Lindelof Existence}
\begin{theorem}
	[Picard-Lindelof Theorem]
	Let \( D \subset \mathbb{R}\times \mathbb{R}^n\) be open, and \( (x_0,y_0) \in D\), \( f:D \to \mathbb{R}^n \). Assume:
	\begin{enumerate}
		\item \( f \) is continuous in the first arguement
		\item \(f\) is Lipschitz in the second variable by constant \( K \) on \( D \)
	\end{enumerate}
	Then there exists \( a>0 \) such that the IVP \[
		\begin{cases}
			y^\prime(x)=f(x,y(x)) \\
			y(x_0)=y_0
		\end{cases}
	\]
	has a solution on \( [x_0-a, x_0+a] \). Further, the solution is unique on that interval. That is to say, any solution defined on the interval must agree with \( y \) for the entire interval.
\end{theorem}
\subsubsection{Proof via Picard Iteration}
\subsubsection{Proof via Banach's Contraction Mapping Theorem}

\begin{proof}
	Choose a closed rectangle on the domain of \( f \):
	\[
		R^\prime=[x_0-A, x_0+A]\times [y_0-L, y_0+L]
	\]
	Recall that \( f \) is continuous, thus it is uniformly continuous and bounded on the compact set \( R^\prime \). Let \[
		|f(x,y)|\leq M
	\] for any \( (x,y)\in R^\prime \). Define length \[
		a\leq A, \quad a\leq \frac{ L }{ M } , \quad aK<1
	\]
	Then let \( I=[x_0-a,x_0+a] \), we have the complete metric space
	\[
		X=\{ y\in C(I, \mathbb{R}^n); |y(x_0)-y_0|\leq L\}
	\]
	, complete as it is closed. We want to show that the picard map \[
		T(y)=y_0+\int_{x_0}^{x} f(s,y(s)) \dd s
	\]
	is a strict contranction from \( X \) to \( X  \), and thus there is a unique solution on the interval \( I \). \\
	Let \( y\in X \), then \[
		\norm{T(y)(x)-y_0}\leq \int_{x_0}^{x} f(s,y(s)) \dd s \leq aM \leq L
	\]
	thus \( Ty\in X \), thus \( T:X\to X \).\\

	\begin{align*}
		\norm{T(y_1)(t)-T(y_2)(t)} & =\norm{\int_{x_0}^{x} f(s, y_1(s))-f(s, y_2(s)) \dd s} \\
		                           & \leq |x-x_0| K||x_1-x_2||_\infty
	\end{align*}
	thus \[
		\norm{T(y_1)-T(y_2)} _\infty \leq Ka \norm{y_1-y_2} _\infty
	\]
	and recall that \( 0<aK<1 \), \( T \) is a strict contraction on \( X \). Then, there exists a unique fixed point of \( T \) the picard map on \( X \), and thus the fixed point \( y^* \) is the unique solution for the IVP.
\end{proof}
\section{Uniqueness Theorems}
\subsection{Picard-Lindelof Uniqueness}
\begin{theorem}
	[Picard-Lindelof Theorem]
	Let \( D \subset \mathbb{R}\times \mathbb{R}^n\) be open, and \( (x_0,y_0) \in D\), \( f:D \to \mathbb{R}^n \). Assume:
	\begin{enumerate}
		\item \( f \) is continuous in the first arguement
		\item \(f\) is Lipschitz in the second variable by constant \( K \) on \( D \)
	\end{enumerate}
	Then there exists \( a>0 \) such that the IVP \[
		\begin{cases}
			y^\prime(x)=f(x,y(x)) \\
			y(x_0)=y_0
		\end{cases}
	\]
	has a solution on \( [x_0-a, x_0+a] \). Further, the solution is unique on that interval. That is to say, any solution defined on the interval must agree with \( y \) for the entire interval.
\end{theorem}
\subsubsection{Proof I, Iterated Picard Style Estimates}
\subsubsection{Proof II, Gronwall's Inequality}
\subsection{Osgood's Uniqueness Theorem}
\section{Gronwall's Inequality}
\begin{theorem}
	[Gronwall's Inequality Integral Form]
	Let \( F:[a,b]\to \mathbb{R} \) and \( G: [a,b]\to [0,\infty)  \) be continuous. Suppose \( y:[a,b]\to \mathbb{R} \) continuous and satisfies \[
		y(t)\leq F(t)+\int_{a}^{t} G(s)y(s) \dd s, \quad \forall t\in[a,b]
	\]
	Then for every \( t\in[a,b] \), \[
		y(t)\leq F(t)+\int_{a}^{t} F(s)G(s)e^{\int_{s}^{t} G(u) \dd u} \dd s
	\]

\end{theorem}
\begin{theorem}
	[Differential Gronwall's Inequality]
	Let \( u, K:[a,b]\to \mathbb{R} \) be respectively differentiable and continuous, and \[
		u^\prime(t)\leq K(t) u(t) , \quad \forall t\in[a,b]
	\]
	then
	\[
		u(t)\leq u(a)e^{\int_{a}^{t} K(s) \dd s}
	\]
\end{theorem}
\subsection{Proof 1}
\subsection{Proof 2}
\subsection{Proof 3}
\subsection{Proof 4}
\section{Maximal Solutions}
\begin{definition}
	We call \( f:A\to \mathbb{R}^n \) locally lipschitz if for every compact subset \( U \) of \( A \) it is lipschitz.
\end{definition}
\begin{theorem}
	Let \( D \subset \mathbb{R}^n \) be a open and connected domain and \( F:D\to \mathbb{R}^n \) locally lipschitz. Then there exists a unique maximal solution \( x:I_{\mathrm{max}} \to D  \) of \[
		x^\prime= F(x(t)),
		x(0)=x_0\in I_{\mathrm{max}}
	\]
	where on a unique interval of existence \( I_{\mathrm{max}}=(a,b)  \) where \( 0\in I_{\mathrm{max}}  \). In addition, if \( b<\infty \), \[
		\lim_{x \to b}|x(t)|+\frac{ 1 }{ d(x(t),\partial D) }
	\]
\end{theorem}

\section{Continuous Dependence}
\begin{theorem}
	[Continuous Dependence on Initial Conditions and Vector Fields]
	Let \( f: [a,b]\times \mathbb{R}^n \to \mathbb{R}^n\), be Lipschitz in the Second variable with constant \( L\geq 0 \).\\

	Let \( x_1, x_2:[a,b]\to \mathbb{R}^n \),  assume \[
		\norm{x_1^\prime(t)-f(t,x_1(t))} <\varepsilon_1, \quad \norm{x_2^\prime(t)-f(t, x_2(t))}<\varepsilon_2
	\]
	and initial values satisfy \[
		\norm{x_1(a)-x_2(a)}< \delta
	\]
	then for all \( t\in[a,b] \),
	\[
		\norm{x_1(t)-x_2(t)} < \delta e^{L(t-a)} + \frac{ \varepsilon_1+\varepsilon_2 }{ L } (e^{L(x-a)} -1)
	\]
\end{theorem}
\section{Autonomous ODE and Flow}

\chapter{Planar Linear Systems of ODE's}


\section{Distinct Real Eigenvalues for a Planar System}
Since a \( 2x2 \) matrix has a two distinct real eigenvalues, we know that we diagonalize the matrix \( A \) into
\[
	D=CAC^{-1}=
	\begin{bmatrix}
		\lambda_1 & 0         \\
		0         & \lambda_2
	\end{bmatrix}
\]
which defines a eigen basis by \( C \), of \( v_1 \) and \( v_2 \).
Thus, it is clear that such an ODE \[
	X^\prime=AX(t)
\]
must satisfy
\[
	x(t)=a(t)v_1+b(t)v_2, \quad x^\prime(t)=a(t)\lambda_1v_1+b(t)\lambda_2v_2
\]
which suggests the solution
\[
	x(t)=ae^{\lambda_1 t} v_1 + b e^{\lambda_2t} v_2
\]
This defines two lines by setting \( a=0 \) or \( b=0 \), and the two directions
\begin{enumerate}
	\item If \( 0<\lambda_1<\lambda_2 \), then we see that along \( v_1 \) and \( v_2 \), the arrows points away from the origin to infinity.
	\item If one is positive one is negative, arrows diverges from origin on the positive line, converges to origin on negative line.
	\item If both are negative, both points to origin
\end{enumerate}
\section{Complex Eigenvalues for a Planar System}
If the Eigenvalues are non-real, then we still have

\section{Repeated Eigenvalues on a Planar System}
Consider that you are given a homogeneous, autonomous, linear ODE:
\[
	x^\prime (t)=A x(t), \quad A\in \mathbb{R}^{2\times 2}
\]
we seek a general solution to this, assuming the solution exist. But by our computation, it has the eigenvalue \( \lambda \) with algebrain multiplicity of \( 2 \) but geometric multiplicity of \( 1 \). We can no longer compute two linear independent eigenvectors to construct our solutions. Thus, we resort to the general eigenspace decomposition. Specifically, we can use Jordan Decomposition (for any dimentions). \\

In this case, we know there is a jordan decomposition of \( A \) into one jordan block, of the form \[
	A=C J C^{-1}
\]
where \[
	J= J(\lambda,2)=
	\begin{bmatrix}
		\lambda & 1       \\
		0       & \lambda
	\end{bmatrix}
\]
where \( v= C v^\prime C^{-1}  \) is an eigenvector of \( A \) where \( v^\prime \) is an eigenvector of \( J \). Then to find the other basis vector for \( J \) we define \[
	(A+\lambda I)u=v
\]
as such a vector is guarenteed to exists by the jordan form. And thus \[
	Au=v+\lambda u
\]

Using \( v \) and \( u \) as a basis will simplify our later calculations.\\

We seek a solution of the form \[
	x(t)=a(t)v+b(t)u
\]
Since \[
	Ax(t)=a(t)Av+b(t)Au=(\lambda a(t)+b(t))v+\lambda b(t)u
\]
we have that \[
	a^\prime(t)=\lambda a(t)+b(t), \quad b^\prime(t)=\lambda b(t)
\]
then \[
	b(t)=Be^{\lambda t}
\]
Since \[
	(e^{-\lambda t} a(t))^\prime=b(t) e^{-\lambda t}=B, \quad a(t)= (A+tB)e^{\lambda t}
\]
Then we have the final form of the solution:
\[
	x(t)=(A+tB)e^{\lambda t}v +Be^{\lambda t} u=e^{\lambda t}(Av+B(tv+u))
\]
\subsection{Qualitative Behaviors}
\begin{enumerate}
	\item If \( \lambda>0  \)
	      \begin{enumerate}
		      \item If \( c_2=0 \) then the solution stays on \( \operatorname{span} v \), moving away form the origin.
		      \item If \( c_2 \neq 0 \), then for large \( t \), \( x(t) \) is approximately \( Be^{\lambda t}(tv+u) \) at time \( t \), moving away from the origin.
	      \end{enumerate}
	\item If \( \lambda<0 \), then the same thing above but towards the origin.
\end{enumerate}
\section{Determinant of the Matrix Exponential}
\begin{theorem}
	[Determinant of a Matrix Exponential]
	If \( B \) is a \( n\times n \) matrix over \( C \) \[
		\operatorname{det}e^B=e^{\operatorname{tr} B}
	\]
\end{theorem}
\begin{proof}
	Recall the definition of \( e^B \):
	\[
		e^B=\sum_{i=0}^{\infty} \frac{ B^i }{ i! }
	\]
	We know that since \( B \) is a complex matrix, it has a jordan decomposition, providing
	\[
		B=CJC^{-1}
	\]
	where \( J \) is a diagonal matrix of square matrices:
	\[ J(\lambda,k) \]
	where \( k \) is the size of the block. We treat the case of a single block:
	\[
		J_i=\lambda I+N
	\]
	where \( N \) is nilpotent. Recall that \[
		\operatorname{det} e^{J_i} =\operatorname{det}e^{\lambda I}\operatorname{det}e^N
	\]
	See that $e^N=\sum_{i=0}^{k} \frac{ N^i }{ i! } $ and since sums and products of a strictly upper triangular matrix \(  N \) is also strictly triangular, the diagonal of \( e^N \) is only contributed by \( I \). Thus, \( e^n \) is upper triangular and \( \operatorname{det} e^N = \Pi_{i=1}^{n} 1 =1\)
	Thus \[
		\operatorname{det}e^{J_i}=\operatorname{det}e^{\lambda I}= \operatorname{det} ((\sum_{i=0}^{\infty} \frac{ \lambda^i }{ i! } )I)=\operatorname{det} e^\lambda = e^{\sum_{i=1}^{k} \lambda} = e^{\operatorname{tr} J(\lambda,k)}
	\]
	The consider all the blocks, with \( n \) dimentional matrix \( B \) and jordan blocks \( J_i \)
	\[
		\operatorname{det}e^B=\operatorname{det}e^J= \Pi_i \operatorname{det} J_i  =e^{\operatorname{tr} B}
	\]
	where we used that \( e^J \) is a identically partitioned block diagonal matrix with respect to \( J \) with the diagonal elements \( e^{J_i} \).

\end{proof}

\begin{theorem}
	[Liouville's formula]
	Consider the Linear System \[
		x^\prime(t)=A(t)x(t), \quad A(t)\in C([t_0,t_1], \mathbb{R}^{n\times n} )
	\]
	A fundamental matrix is a matrix-valued solution \( \Phi(t) \) of
	\[
		\Phi^\prime(t)=A(t)\Phi(t), \quad \Phi(t_0)=B_0\in \operatorname{GL}(n, \mathbb{R})
	\]
	Then \[
		\operatorname{det}\Phi(t)= \operatorname{det} (B_0) e^{\int_{t_0}^{t} \operatorname{tr} (A(s))\dd s}
	\]
	Equivalently, the Wrongskian \( W(t) = \operatorname{det}\Phi\) satisfies the scalar ODE \[
		W^\prime(t)= \operatorname{tr} (A(t)) W(t), \quad W(t_0)=\operatorname{det}B_0
	\]
	(Assuming the solution exist and is unique)
\end{theorem}
\begin{proof}
	We have that \[
		(\Phi(t)^\prime -A(t)\Phi(t)) e^{-\int_{t_0}^{t} A(s) \dd s }=0
	\]
	Thus \[
		(\Phi(t)e^{-\int_{t_0}^{t} A(s)\dd s}) ^\prime=0
	\]
	and \[
		\Phi(t)=B_0 e^{\int_{t_0}^{t} A(s) \dd s}
	\]
	and thus \[
		\operatorname{det}\Phi(t)= \operatorname{det}B_0 e^{\operatorname{tr} \int_{t_0}^{t} A(s) \dd s} = \operatorname{det}(B_0) e^{\int_{t_0}^{t} \operatorname{tr} A(s) \dd s}
	\]
	Then we have shown the first part. Now let \( W(t)= \operatorname{det}\Phi(t) \). \[
		W^\prime(t)=\operatorname{det}(B_0) \operatorname{tr} (A(t)) e^{\int_{t_0}^{t} \operatorname{tr} A(s) \dd s } = \operatorname{tr} (A(t)) W(t)
	\]
\end{proof}
\section{General Linear First Order Systems}
Let us first define a few preliminary concepts.\\

We call a matrix \( A(t) \) differentiable if all of its coefficients are. In this case, we denote for \[
	(A(t))_{jk} =a_{jk} (t)
\], that
\[
	(A^\prime(t))_{jk} = a^\prime_{jk} (t)
\]
Thus several common rules of differentiation follows
\begin{enumerate}
	\item Product rule \[
		      ((AB)^\prime(t))_{jk} =(\sum_{i=1}^{n} (a_{ji}(t)b_{ik} (t))^\prime)_{jk}=(A^\prime(t)B(t)+A(t)B^\prime(t))_{jk}
	      \]
	\item If \( \det A(t) \neq 0 \), \[
		      DA(t)^{-1} =-A(t)^{-1} \dot{A}(t)A(t)^{-1}
	      \]
\end{enumerate}
\begin{theorem}
	The linear first-order system \[
		\dot{x}(t)=A(t)x(t)
	\]
	where \( A(t)\in C(I, \mathbb{R}^{n\times n} ) \) where \( I\subset \mathbb{R} \) is some interval. Then there is a unique solution satisfying the initial condition \( x(t_0)=x_0\in \mathbb{R}^{n}\) defined for all \( t\in I \).
	\label{thm:Linear First Order System}
\end{theorem}
\begin{proof}
	We do this proof as a simple practice.\\

	First, we wish to apply Picard-Lindelof locally. Since \( f(t)=A(t)x(t) \) is \( C^1 \), then there is a costant \( \alpha \) such that there exists a solution \( x:[t_0-\alpha, t_0+\alpha ]\to \mathbb{R}^n \) unique on its interval of definition. We may extend it by defining new IVP with \( x(t_{\pm 1} )=x(t_0\pm \alpha ) \), glueing gives the unique solution on the union of the defining intervals of the solution. By repeated application of this process, we must come to a maximal interval \( I_{\max} \subset I\) (of course there must be a supremum defined by the total order on the intervals, and the supremum being \( I \)). \\

	Consider then case where \( I_{\max} \subsetneq I\). Suppose for the sake of contradiction that \( I_{\max} \) is closed. Then, \( I_{\max}=[a,b] \), but we may still extend it using Picard-Lindelof, thus \( I_{\max} \) must be open. Write \( I_{\max} = (I_-,I_+) \). If \( I_+ (I_-)\in \operatorname{int} I   \), then define an open ball centered at \( I_+ (I_-)  \), we may extend the solution again. Thus by contradiction, if \( I \)  is open then \( I=I_{\max} \). Otherwise, if \( I \) is not open, we have that for a side not open, (suppose the left side)
	\begin{align*}
		x(t)                      & =x_0+\int_{t_0}^{t} A(s)x(s)\mathrm{d}s             \\
		x(a)                      & =x_0+\int_{t_0}^{a} A(s)x(s)\mathrm{d}s             \\
		\frac{ x(a)-x(t) }{ a-t } & =\frac{ 1 }{ a-t } \int_{t}^{a} A(s)x(s)\mathrm{d}s
	\end{align*}
	Taking the limit of both sides we have that \[
		x^\prime(a)=A(a)x(a)
	\]
	where the derivative is defined on a single side.
	Applying to the other side the same way, we find that \( I_{\max} \) may still be extended. Thus, \( I=I_{\max} \) is the maximal interval of existence of a unique solution.
\end{proof}
It is a simple observation that
\begin{enumerate}
	\item \( C(I,\mathbb{R}^n)\) is a vector space
	\item \( 0 \) is a solution to the system in theorem \ref{thm:Linear First Order System}
	\item Linear combinations of solutions are also solutions
\end{enumerate}
Thus, the solutions forms a subspace of \( C(I,\mathbb{R}) \). Further, a solution defined by \( \phi(t,x_0) \) can be written as \[
	\phi(t, t_0 ,x_0)=\sum_{i=1}^{n} \phi(t,t_0,e_i) (x_0)_i
\]
where \( e_i \) is the \( i \)th standard basis. Thus, define
\begin{equation}
	\Pi(t)=
	\begin{bmatrix}
		\phi(t,t_0, e_1) & \cdots & \phi(t,t_0, e_n)
	\end{bmatrix}
	\label{eq:Fundamental Matrix}
\end{equation}
See that this solve the first order system \[
	\Pi^\prime(t)=A(t)\Pi(t), \quad \Pi(t_0)=\mathbb{I}
\]
and has the property \[
	\Pi(t)x_0=\phi(t,t_0,x_0)
\]
Thus we define
\begin{definition}
	The matrix defined by \eqref{eq:Fundamental Matrix} is the \textbf{Principal Matrix Solution} (at \( t_0 \)), solving the \textbf{Matrix Valued} IVP \[
		\Pi(t,t_0)=A(t)\Pi(t), \qquad \Pi(t_0)=\mathbb{I}
	\]
\end{definition}
Here is an interesting property of the Principal Matrix Solution:\[
	\Pi(t,t_0)=\Pi(t,t_1)\Pi(t_1,t_0)
\]
Let \( t=t_1 \), we obtain
\[
	\mathbb{I}=\Pi(t_0,t_1)\Pi(t_1,t_0)
\]
Thus,
\[
	(\Pi(t,t_0))^{-1} =\Pi(t_0,t)
\]
Take this to a more general perspective, take \( n \) arbitrary solutions \( \{ \phi_1, \cdots, \phi_n \}\),
we an define a matrix valued solution \[
	U=[\phi_1 \quad \cdots \quad \phi_n ]
\]
Then we have that \[
	\dot{U}(t)U^{-1}(t) =A(t)
\]
\begin{definition}
	Define \[
		U(t)=
		\begin{bmatrix}
			\phi_1 & \cdots & \phi_n
		\end{bmatrix}
	\]
	where \( \phi_i \) are solutions described by theorem \ref{thm:Linear First Order System}. The determinant of \( U \) is called the \textbf{Wronski Determinant}
	\begin{equation*}
		W(t)=\det(\phi_1, \cdots, \phi_n)
	\end{equation*}
	In addition, if \( W(t)\neq 0 \), \( U \) is called a fundamental matrix solution.
\end{definition}
\begin{proposition}
	The Wronski determinant of \( n \) solutions satisfies
	\begin{equation*}
		W(t)=W(t_0)\int_{t_0}^{t} \tr(A(s))\mathrm{d}s
	\end{equation*}
\end{proposition}
\begin{proof}
	Proved in class
\end{proof}
\subsection{Inhomogeneous Systems}
We define the \( n \) dimentional Linear System:
\begin{equation}
	\dot{x}=A(t)x(t)+f(t)
	\label{eq: n dimentional Inhomogeneous Linear System}
\end{equation}
where \( A(t)\in C(I,\mathbb{R}^{n\times n} ) \) and \( f(t)\in \mathbb{R}^n \). See that if \( x,y \) are solutions to the system \eqref{eq: n dimentional Inhomogeneous Linear System}, their difference must take the form \[
	(x-y)^\prime(t)=A(t)(x(t)-y(t))
\]
a solution to the homogeneous system. With this as motivation, we can easily find a solution to the inhomogeneous system.
\begin{theorem}
	The solution for the \( n \) dimentional Inhomogeneous linear IVP \[
		x^\prime(t)=A(t)x(t)+f(t), \quad x(t_0)=x_0
	\]
	where \( A(t)\in C(I,\mathbb{R}^{n\times n}  \) and \( f\in C(I,\mathbb{R}^n) \) is given by
	\begin{equation}
		x(t)=\Pi(t,t_0)x_0+\int_{t_0}^{t} \Pi(s,t_0)f(s)\mathrm{d}s
		\label{eq:Solution for n dimentional Inhomogeneous Linear IVP}
	\end{equation}
	\label{thm: solution for the n dimentional Inhomogeneous IVP}
\end{theorem}
\chapter{Dynamical Systems}
We will first provide an abstract definition of a dynamical system, and illustrate with examples later.
\begin{definition}
	A dynamical system is a semigroup \( (G,\cdot) \) with identity \( e \) acting on a set \( M \).
\end{definition}
Then, the action of the semigroup is given by \[
	T:G\times M\to M
\]
with \[
	T(g,x)=T_g(x)
\]
following \[
	T_g\circ T_h=T_{g\cdot h}
\]
In addition, if \( (G,\circ) \) is a group, we call dynamical system invertable. We will mostly discuss the following discrete and continous (semi)-groups:
\[
	(\mathbb{R}, +) \quad (\mathbb{R}^+, +), \quad (\mathbb{Z}, +), \quad (\mathbb{N}^0, +)
\]
As a elementary example
\begin{example}
	With \( (\mathbb{N}^0, +) \) acting on interval \( I\subset \mathbb{R} \), define the action with the function $f:I\to I$, \[
		T(n,x)=f_n(x), \quad f_n=f\circ f_{n-1} \text{ and } f_0(x)=(x)
	\]
	Of course, if \( f \) is invertible and defined by \( (\mathbb{Z},+) \), this would then be an invertable dynamic system.
\end{example}
Or a more relevant example to our study
\begin{example}
	Take group \( (\mathbb{R}, +) \) and let \( M= \mathbb{R}^n \), we can define the flow of an autonomous linear differential equation \[
		\Phi
	\]
	with the action
	\[
		\Phi(t,x) = \phi_t(x)
	\]
	with the property that \( x(t)=\phi_t(x_0) \) solves
	\[
		x^\prime=Ax, \quad x(0)=x_0
	\] with \( A\in \mathbb{R}^{n\times n}   \)for all \( x_0\in \mathbb{R} \). We will expand on this in the following section.
\end{example}
\section{Flow of an Autonomous Equation}
First we present the general IVP of the Autonomous problem
\begin{equation}
	x^\prime=f(x),\quad x(0)=x_0
	\label{eq:autonomous IVP for flow}
\end{equation}
where we assume \( f\in C^k(M, \mathbb{R}^n) \) with \( k\geq 1 \) and \( M\subset \mathbb{R}^n \) open. Such a system (specifically \( f \)) can be considered a vector field on \( M \), and solutions to the IVP must be tangent to this vector field at each point. The parametric curve defined by the solution is also called an integral curve. Specifically, a solution \( x \) is called a integral curve on \( x_0 \) if \( x(0)=x_0 \). \\

For every point \( x_0\in M \), we know that there exists a unique maximal solution \( x(t) \) (maximal integral curve) to the IVP with \( x(0)=x_0 \), on a maximal interval \( I_{x_0} =(T_-(x_0), T_+(x_0)) \). Define the set \[
	W= \bigcup_{x \in M} I_{x}\times \{x\} \subset \mathbb{R}\times M
\]
Then we may define the flow of the differential equation \[
	\Phi:W\to M, \quad (t,x_0) \mapsto \phi(t,x_0)
\]
Where \( \phi(t,x_0) \) is the maximal integral curve at \( x_0 \). Let \( x(t)=\phi(t,x_0) \) define maximan integral curve at \( x_0 \), then $y(t)=\phi(t+s,x_0)$ with \( y(0)=x(s) \) must define the maximal integral curve at \( x(s) \). Their maximal intervals of existence necessarily statisfy \( I_{x_0} =s+ I_{x_s} \). As an interesting conclusion, \begin{equation}
	\phi(t+s,x_0)=\phi(t, \phi(s,x_0))
	\label{eq: semigroup action}
\end{equation}
In the case where the ODE is linear (and some other cases), clearly \( I_{x_0} = \mathbb{R} \) for each \( x_0 \) meaning \( W=\mathbb{R}\times M \), thus the flow under this assumption is a dynamical system defined by \( (\mathbb{R}, +) \) on \( M \) with the action \[
	\Phi:\mathbb{R}\times M \to M, \quad (t,x_0)\to \phi(t,x_0)
\]
We can verify that this indeed follows the semigroup action: write \( \phi(t,x)=\phi_t(x) \), then \[
	(\phi_a\circ \phi_b ) (x) = \phi_{a+b} (x)
\] which follows directly from \eqref{eq: semigroup action}.\\

Hence, we have 3 different view points of the map \( \Phi \):
\begin{enumerate}
	\item The first level is simply the master equation of \[
	 \Phi:W \times M
	\]
\item Fix \( x\in M \), we obtain a curve of the solution to the ODE parametrizing with respect to \( t \):\[
 \phi_x:I_x \to M
\]
\item Fix \( t \) in a suitable subset \( \overline{I}\subset \bigcup I_x\) and a corresponding subset \( X\subset M \), we have a operator on \( X \) defined \[
 \phi_t: X\to X 
\]
where the output of the operator highlights the dynamics of the system depending on the initial value \( x \).
\end{enumerate}
\begin{theorem}
    Let \( f\in C^k(M, \mathbb{R}^n) \) , then for each \( x\in M \), there is a unique pair maximal integral curve \[
     \phi_x: I_{x}\to M, \quad \phi_x(0)=x 
    \]
	Define 
	\[
		W=\{(t,x)\in \mathbb{R}\times M: t\in I_x\}, \quad \Phi:W\to M, \quad \Phi(t,x)=\phi_x(t)
	\]
	we have that \( M \) is open and \( \Phi\in C^k(W,M) \) a (local) flow on \( M \), that is, it satisfies \[
	 \Phi(t+s,x)=\Phi(t,\Phi(s,x)), \quad \Phi(0,x)=x
	\]
\end{theorem}
In the operator viewpoint, set \( t_1, -t_1\in I \) and thus we obtain an inverse to the operator.  
Taking it a step further, we can formalize all these concept algebraicly
\section{Exursion: Global Flow}
\begin{definition}
    A \textbf{Lie Group} is a group that is also a smooth manifold
\end{definition}
\begin{definition}
 Let \( I \) be a Lie Group with an identity element \( e \), \( M \) be a smooth manifold. If \[
  \Phi: I\times M\to M
 \]
 is a smooth map satisfying \[
  \Phi(e,x)=x
 \]
 and
 \[
  \Phi(g,\Phi(h,p))=\Phi(gh,p)
 \]
 We call \( \Phi \) a smooth action of \( I \) on \( M \). For each \( g\in I \) we define \[
  \Phi_g:M \to M, \quad \Phi_g(p)=\Phi(g,p)
 \]
\end{definition}
\begin{theorem}
 If \( \Phi  \) is a smooth action of Lie Group \( G \) on smooth manifold \( M \), then for every \( g\in G \), the map \( \Phi_g \) is a diffeomorphism of \( M \). Moreover, \[
  \Phi_{gh} = \Phi_g \circ \Phi_h, \quad \Phi_e=\mathbb{I}_M, \quad \Phi_{g^{-1} } =(\Phi_g)^{-1} 
 \]
 so the action defined a group homeomorphism \[
  G \to \operatorname{Diff}(M), \quad g\mapsto \Phi_g 
 \]
\end{theorem}
\begin{definition}
 A family $\left\{\Phi_t\right\}_{t \in \mathbb{R}} \subset \operatorname{Diff}(M)$ is called a one-parameter transformation group on $M$ if:\\
1. $\Phi_0=\mathrm{id}_M$,\\
2. $\Phi_{t+s}=\Phi_t \circ \Phi_s$ for all $s, t \in \mathbb{R}_1$\\
3. the map
$$
\Phi: \mathbb{R} \times M \rightarrow M, \quad(t, p) \mapsto \Phi_t(p)
$$
is smooth.\\
\end{definition}


Thus a one-parameter transformation group is exactly a smooth action of the Lie group $(\mathbb{R},+)$ on $M$. 
\begin{definition}
 
\end{definition}

\end{document}
